\chapter{Noções de Direitos Humanos}

\section{Conceito e fundamentos dos direitos humanos}

Direitos humanos são posições jurídicas destinadas a proteger a dignidade, a liberdade, a igualdade e as condições materiais de existência das pessoas. No plano internacional, constituem um sistema de normas, princípios, instituições e mecanismos de proteção. No plano interno, muitos desses direitos são reconhecidos como direitos fundamentais pela Constituição.

A distinção entre \termo{direitos humanos} e \termo{direitos fundamentais} é didática. Em geral, a primeira expressão enfatiza a proteção internacional e universal; a segunda, a positivação constitucional dentro de determinado Estado. Os dois planos se comunicam continuamente.

\subsection{Dignidade da pessoa humana}

Dignidade não é apenas um princípio abstrato. Ela funciona como fundamento para reconhecer autonomia, integridade física e psíquica, igualdade, condições mínimas de existência e proibição de tratamento degradante.

\section{Características dos direitos humanos}

A doutrina utiliza características que ajudam a compreender o regime protetivo.

\subsection{Universalidade}

Os direitos humanos pertencem a todas as pessoas, sem depender de nacionalidade, raça, religião, sexo ou outra condição. Universalidade não significa que a concretização seja idêntica em todos os contextos jurídicos, mas afasta exclusões arbitrárias do status de sujeito de direitos.

\subsection{Historicidade}

A formulação e a proteção de direitos desenvolvem-se historicamente. Novas demandas sociais podem ampliar formas de proteção sem eliminar direitos anteriormente reconhecidos.

\subsection{Indivisibilidade e interdependência}

Direitos civis, políticos, econômicos, sociais e culturais não formam compartimentos estanques. Liberdade de expressão pode depender de acesso à educação; participação política pode ser esvaziada por exclusão social extrema.

\subsection{Inalienabilidade e irrenunciabilidade}

Direitos ligados à própria dignidade não podem ser tratados simplesmente como mercadoria. A possibilidade de exercício ou de não exercício de determinada faculdade não equivale a autorizar supressão absoluta do direito.

\subsection{Relatividade}

A maioria dos direitos não é absoluta. Conflitos entre direitos podem exigir harmonização segundo legalidade, finalidade legítima e proporcionalidade. A proibição de tortura, entretanto, possui regime internacional especialmente absoluto.

\section{Dimensões de direitos}

A classificação em dimensões é uma ferramenta histórica, não uma cronologia de substituição.

\begin{table}[ht]
\centering
\small
\begin{tabularx}{\textwidth}{l X X}
\toprule
Dimensão & Ênfase & Exemplos \\
\midrule
Primeira & liberdade e limitação do poder & vida, liberdade, propriedade, devido processo, participação política. \\
Segunda & igualdade material e prestações & saúde, educação, trabalho, previdência e assistência. \\
Terceira & solidariedade e interesses coletivos/difusos & meio ambiente, desenvolvimento, paz e patrimônio comum. \\
\bottomrule
\end{tabularx}
\end{table}

A expressão ``gerações'' pode induzir à ideia errada de que uma etapa substitui outra. Os direitos acumulam-se e interagem.

\section{Sistema global de proteção}

Após a Segunda Guerra Mundial, a proteção internacional passou por forte institucionalização no âmbito das Nações Unidas.

\subsection{Declaração Universal dos Direitos Humanos}

A Declaração Universal dos Direitos Humanos (DUDH) foi adotada pela Assembleia Geral da ONU em 10 de dezembro de 1948. Ela articula direitos civis, políticos, econômicos, sociais e culturais sob o fundamento da dignidade e da igualdade.

Entre seus conteúdos estão:
\begin{itemize}
\item igualdade e não discriminação;
\item vida, liberdade e segurança;
\item proibição de escravidão e tortura;
\item reconhecimento perante a lei;
\item garantias contra prisão arbitrária;
\item julgamento justo;
\item privacidade e família;
\item liberdade de circulação, pensamento, consciência, religião e expressão;
\item reunião e associação;
\item participação política;
\item trabalho, proteção social e educação;
\item participação na vida cultural.
\end{itemize}

A DUDH é uma declaração da Assembleia Geral, não um tratado multilateral ratificado segundo o procedimento convencional. Isso não reduz sua importância histórica, jurídica e interpretativa.

\section{Pactos internacionais de 1966}

A DUDH foi posteriormente complementada por tratados juridicamente vinculantes.

\subsection{Pacto Internacional sobre Direitos Civis e Políticos}

O PIDCP protege, entre outros, direito à vida, liberdade, integridade, devido processo, liberdade religiosa, expressão, reunião, associação e participação política.

O pacto prevê deveres imediatos de respeito e garantia e estabelece mecanismo internacional de supervisão pelo Comitê de Direitos Humanos.

\subsection{Pacto Internacional sobre Direitos Econômicos, Sociais e Culturais}

O PIDESC protege direitos como trabalho, condições justas, sindicalização, seguridade social, padrão de vida adequado, saúde, educação e participação cultural.

Sua implementação combina obrigações imediatas, como não discriminação, e realização progressiva mediante utilização do máximo de recursos disponíveis, nos termos do tratado.

\begin{teoria}[Carta Internacional de Direitos Humanos]
A DUDH, o PIDCP e o PIDESC formam o núcleo frequentemente chamado de \textbf{International Bill of Human Rights}. Os pactos transformaram muitos princípios da Declaração em obrigações convencionais específicas.
\end{teoria}

\section{Proteção internacional e Constituição brasileira}

A Constituição de 1988 abre o sistema constitucional à proteção internacional dos direitos humanos.

\subsection{Tratados de direitos humanos}

O art. 5º, §3º, prevê que tratados e convenções internacionais sobre direitos humanos aprovados em cada Casa do Congresso, em dois turnos, por três quintos dos votos, equivalem a emendas constitucionais.

Para tratados de direitos humanos incorporados sem esse rito qualificado, a jurisprudência do STF reconhece status supralegal, acima da legislação ordinária e abaixo da Constituição.

Essa distinção é importante para compreender controle de convencionalidade e conflitos entre normas internas e obrigações internacionais.

\section{Igualdade e não discriminação}

Igualdade possui dimensão \termo{formal} e \termo{material}. A igualdade formal impede discriminações arbitrárias perante a lei. A igualdade material reconhece que situações historicamente desiguais podem exigir medidas diferenciadas para produzir acesso efetivo a direitos.

\subsection{Discriminação direta e indireta}

Discriminação direta ocorre quando a distinção utiliza explicitamente critério proibido ou injustificado. Discriminação indireta pode ocorrer quando regra aparentemente neutra produz desvantagem desproporcional sobre determinado grupo sem justificativa adequada.

\begin{exemplo}[barreira indireta]
Um serviço público exclusivamente digital exige interação por interface incompatível com leitor de tela. A regra é igual para todos, mas o desenho tecnológico impede pessoas cegas de exercer o direito em igualdade de condições. A ausência de distinção formal não elimina o efeito discriminatório.
\end{exemplo}

\section{Agenda 2030 e Objetivos de Desenvolvimento Sustentável}

A Agenda 2030, adotada em 2015 pelos Estados-membros da ONU, estabelece 17 Objetivos de Desenvolvimento Sustentável e 169 metas.

Os ODS combinam dimensões econômica, social, ambiental e institucional. Para controle público, alguns vínculos são especialmente relevantes:
\begin{itemize}
\item ODS 5 — igualdade de gênero;
\item ODS 10 — redução das desigualdades;
\item ODS 11 — cidades e comunidades sustentáveis;
\item ODS 16 — paz, justiça e instituições eficazes, responsáveis e inclusivas;
\item ODS 17 — meios de implementação e parcerias.
\end{itemize}

ODS são objetivos de política e cooperação internacional. Não se deve concluir que cada meta da Agenda 2030, isoladamente, crie direito subjetivo judicialmente autoexecutável com conteúdo idêntico ao de uma norma constitucional ou legal.

\section{Convenção sobre os Direitos das Pessoas com Deficiência}

A Convenção da ONU sobre os Direitos das Pessoas com Deficiência foi incorporada ao ordenamento brasileiro com o rito do art. 5º, §3º, da Constituição, adquirindo equivalência constitucional.

A Convenção consolida o \termo{modelo social da deficiência}: a deficiência não é compreendida apenas como condição individual, mas como resultado da interação entre impedimentos e barreiras ambientais, comunicacionais, sociais e atitudinais.

\section{Lei Brasileira de Inclusão — Lei nº 13.146/2015}

A LBI concretiza no plano interno diversos comandos da Convenção.

\subsection{Conceito de pessoa com deficiência}

Pessoa com deficiência é aquela que possui impedimento de longo prazo de natureza física, mental, intelectual ou sensorial que, em interação com uma ou mais barreiras, pode obstruir sua participação plena e efetiva na sociedade em igualdade de condições.

A definição possui dois elementos:
\[
\text{impedimento de longo prazo}
+
\text{interação com barreiras}
\rightarrow
\text{restrição de participação}.
\]

O diagnóstico médico isolado não é suficiente para compreender juridicamente todas as dimensões da deficiência.

\subsection{Avaliação biopsicossocial}

Quando necessária, a avaliação considera:
\begin{itemize}
\item impedimentos nas funções e estruturas do corpo;
\item fatores socioambientais, psicológicos e pessoais;
\item limitação no desempenho de atividades;
\item restrição de participação.
\end{itemize}

\section{Acessibilidade e barreiras}

Acessibilidade é a possibilidade de alcance e utilização, com segurança e autonomia, de espaços, mobiliários, transportes, informação, comunicação, sistemas e tecnologias.

A LBI classifica barreiras em diferentes categorias.

\begin{table}[ht]
\centering
\small
\begin{tabularx}{\textwidth}{l X}
\toprule
Barreira & Exemplo \\
\midrule
Urbanística & calçada sem rota acessível. \\
Arquitetônica & prédio público sem acesso adequado. \\
Transportes & veículo ou estação inacessível. \\
Comunicação e informação & documento sem formato acessível. \\
Atitudinal & comportamento que exclui ou infantiliza a pessoa. \\
Tecnológica & sistema ou aplicativo incompatível com tecnologia assistiva. \\
\bottomrule
\end{tabularx}
\end{table}

\subsection{Acessibilidade digital}

A LBI determina acessibilidade em sítios da Internet mantidos por órgãos de governo e por empresas abrangidas pela regra legal, observadas melhores práticas e diretrizes internacionais.

Acessibilidade digital envolve, por exemplo:
\begin{itemize}
\item navegação por teclado;
\item estrutura semântica interpretável por leitores de tela;
\item contraste adequado;
\item alternativas textuais para conteúdo visual;
\item legendas e recursos para conteúdo multimídia;
\item mensagens de erro compreensíveis;
\item foco visível e ordem lógica de navegação.
\end{itemize}

\section{Adaptação razoável e desenho universal}

\termo{Adaptação razoável} é modificação ou ajuste necessário e adequado que não imponha ônus desproporcional ou indevido, requerido em caso concreto para garantir exercício de direitos em igualdade de condições.

\termo{Desenho universal} busca conceber produtos, ambientes, programas e serviços para serem usados pelo maior número possível de pessoas desde a origem, sem necessidade de adaptação posterior, embora tecnologias assistivas específicas possam continuar necessárias.

A recusa de adaptação razoável integra o conceito de discriminação em razão da deficiência.

\section{Capacidade civil da pessoa com deficiência}

A deficiência não afeta, por si só, a plena capacidade civil. A LBI rompe com modelos que equiparavam automaticamente deficiência intelectual ou mental a incapacidade.

Curatela é medida extraordinária, proporcional às necessidades e circunstâncias de cada caso e restrita aos atos de natureza patrimonial e negocial nos termos legais. Direitos existenciais permanecem protegidos.

\section{Tecnologia assistiva}

Tecnologia assistiva compreende produtos, equipamentos, dispositivos, recursos, metodologias, estratégias, práticas e serviços que objetivem promover funcionalidade, autonomia, independência, qualidade de vida e inclusão.

Leitor de tela, linha Braille, teclado adaptado, prótese e sistemas de comunicação aumentativa são exemplos em contextos distintos.

\section{Necessidades complexas de comunicação — atualização de 2025}

A Lei nº 15.249/2025 alterou a Lei de Acessibilidade e a LBI para incluir a categoria de \termo{pessoa com necessidades complexas de comunicação} e ampliar medidas de comunicação aumentativa e alternativa.

A norma prevê, entre outros pontos, uso de pranchas de baixa tecnologia com pictogramas em espaços públicos e medidas específicas em saúde, educação, cultura e ambientes de uso coletivo, observados os dispositivos legais.

Essa atualização é relevante porque amplia o conceito de acessibilidade comunicacional para além das formas tradicionais de comunicação escrita, oral, Libras ou Braille.

\section{Lei nº 10.098/2000 — Lei de Acessibilidade}

A Lei nº 10.098/2000 estabelece normas gerais e critérios básicos para promoção da acessibilidade de pessoas com deficiência ou mobilidade reduzida.

Ela abrange ambiente urbano, edificações, transportes e comunicação e foi atualizada ao longo do tempo, inclusive pela Lei nº 15.249/2025 quanto a pessoas com necessidades complexas de comunicação.

\section{Atendimento prioritário — Lei nº 10.048/2000}

A redação vigente garante atendimento prioritário, entre outros, a:
\begin{itemize}
\item pessoas com deficiência;
\item pessoas com transtorno do espectro autista;
\item pessoas com 60 anos ou mais;
\item gestantes;
\item lactantes;
\item pessoas com criança de colo;
\item pessoas obesas;
\item pessoas com mobilidade reduzida;
\item doadores de sangue.
\end{itemize}

Os doadores de sangue são atendidos após os demais beneficiários do rol e devem apresentar comprovante de doação com validade de 120 dias, conforme a lei.

\begin{atencao}
Listas antigas de atendimento prioritário podem estar desatualizadas. TEA e doadores de sangue constam expressamente da redação atual.
\end{atencao}

\section{Estatuto da Igualdade Racial — Lei nº 12.288/2010}

O Estatuto busca assegurar à população negra igualdade de oportunidades, defesa de direitos étnicos e combate à discriminação e intolerância étnico-racial.

\subsection{Discriminação racial ou étnico-racial}

Consiste em distinção, exclusão, restrição ou preferência baseada em raça, cor, descendência ou origem nacional ou étnica que tenha por objeto ou resultado restringir o reconhecimento ou exercício de direitos em igualdade de condições.

\subsection{Desigualdade racial}

Refere-se a diferenciação injustificada no acesso e fruição de bens, serviços e oportunidades em razão de raça, cor, descendência ou origem.

\subsection{Ações afirmativas}

São programas e medidas especiais adotados pelo Estado ou iniciativa privada para corrigir desigualdades raciais e promover igualdade de oportunidades.

Ação afirmativa não contradiz automaticamente a igualdade. É instrumento de igualdade material quando fundada em finalidade constitucional e critérios juridicamente adequados.

\section{Políticas públicas e direitos humanos}

Direitos humanos possuem dimensão institucional e orçamentária. Uma política pode existir formalmente e fracassar em alcançar seu público.

Para avaliar efetividade, é útil distinguir:
\begin{itemize}
\item \textbf{disponibilidade}: o serviço existe?;
\item \textbf{acessibilidade}: a população consegue utilizá-lo?;
\item \textbf{aceitabilidade}: respeita dignidade e características dos usuários?;
\item \textbf{qualidade}: atende padrões mínimos adequados?;
\item \textbf{equidade}: grupos vulneráveis enfrentam barreiras desproporcionais?
\end{itemize}

Essa estrutura é especialmente útil para compreender saúde, educação, assistência, acessibilidade e serviços digitais.

\section{Indicadores e desigualdades}

Médias agregadas podem ocultar desigualdades. Um programa pode atingir 90\% da população e ainda excluir quase completamente determinado grupo minoritário.

Por isso, análise de direitos humanos pode exigir dados desagregados por território, raça/cor, sexo, idade, deficiência e outras variáveis juridicamente pertinentes, sempre respeitando proteção de dados e finalidade do tratamento.

\begin{exemplo}[média que oculta exclusão]
Se 95\% dos usuários sem deficiência concluem um serviço digital, mas apenas 40\% dos usuários que dependem de leitor de tela conseguem concluir, a taxa global pode parecer satisfatória e ainda esconder barreira grave de acessibilidade.
\end{exemplo}

\section{Síntese conceitual}

\mapafuncional{Direitos Humanos}
{Dignidade + igualdade + liberdade + condições materiais formam um sistema integrado de proteção}
{Plano internacional $\rightarrow$ DUDH + PIDCP + PIDESC + convenções específicas.\\Plano interno $\rightarrow$ Constituição + tratados + leis.\\PcD $\rightarrow$ impedimento + barreiras + participação.\\Igualdade racial $\rightarrow$ não discriminação + ações afirmativas.\\Políticas públicas $\rightarrow$ acesso + qualidade + equidade.}
{Tratado com rito do art. 5º, §3º: equivalência constitucional.\\PcD: examine barreira, não só diagnóstico.\\Serviço digital: acessibilidade também é requisito de igualdade.\\Média boa: verifique grupos desagregados.}
{Direitos humanos $\neq$ direitos absolutos em geral.\\Deficiência $\neq$ incapacidade civil.\\Igualdade $\neq$ tratamento idêntico sempre.\\Ação afirmativa $\neq$ discriminação arbitrária.\\ODS $\neq$ norma autoexecutável em toda meta.}
{Qual direito está protegido?\\Qual é a fonte normativa?\\Há barreira ou discriminação?\\Existe adaptação razoável?\\Os indicadores revelam desigualdades?\\A medida preserva autonomia e dignidade?}

\begin{resumo}
Direitos humanos devem ser compreendidos como sistema multinível. A proteção internacional influencia a Constituição e a legislação interna; igualdade material explica ações afirmativas e adaptações; e o modelo social da deficiência mostra que exclusão pode estar no ambiente, inclusive em sistemas digitais, e não apenas na condição individual da pessoa.
\end{resumo}
